\chapter{Creio no Espírito Santo}
\chaptermark{Creio no Espírito Santo}

Este capítulo inicia o terceiro artigo do Símbolo Apostólico, professando a fé no Espírito Santo, terceiro Pessoa da Santíssima Trindade. Desdobramos os parágrafos correspondentes do Catecismo da Igreja Católica (CIC 683--747), explorando o dom da fé, a missão conjunta com o Filho, nomes, títulos, símbolos e preparação para Cristo na história da salvação. Para adultos em processo de iniciação cristã pelo RICA, o Espírito desperta a fé batismal, guiando o catecumenato rumo à Vigília Pascal.

\section{O Espírito Santo, Dom da Fé e da Vida Divina (\S 683--716)}

``Ninguém pode dizer `Jesus é Senhor' senão pelo Espírito Santo'' (1Cor 12,3): o Espírito desperta a fé, comunicando intimamente a vida do Pai no Filho pelo Batismo, primeiro sacramento da fé.

Revelado por último na Trindade por ``condescendência divina'', o Espírito habita entre nós, concedendo visão clara após Pai e Filho proclamados.

Cremos no Espírito como Pessoa divina, consubstancial ao Pai e ao Filho, cultuado e glorificado com Eles, no ``economia'' salvífica.

Na missão conjunta Filho-Espírito, distintos mas inseparáveis: o Espírito unge Jesus como Cristo, revela-o invisível, glorifica-o desde a Encarnação até Pentecostes.

Nome, títulos e símbolos revelam-no: ``Santo Espírito'' (ruah, sopro), Paráclito (Consolador, Advogado), Espírito de Verdade, Promessa, Adoção, Cristo; símbolos como água (Batismo), unção (Confirmação), fogo, nuvem/pomba, imposição de mãos.

Para adultos a partir dos 20 anos no RICA, esses mistérios iluminam os exorcismos e uncções do catecumenato quaresmal, invocando o Espírito para purificação e unção batismal na noite pascal.

\section{O Espírito Santo na Preparação Messiânica (\S 717--747)}

O Espírito enche João Batista desde o ventre materno, forerunner de Cristo, cheio do fogo espiritual como Elias, preparando o povo para o Senhor.

``Mais que profeta'', João completa os profetas, voz do Consolador, batizando com água para testemunhar a Luz, vendo o Espírito descer sobre o Cordeiro de Deus.

No Espírito, inicia restauração da ``semelhança divina'', prefigurando Batismo de água e Espírito como novo nascimento.

Maria, obra-prima da missão Filho-Espírito, cheia de graça, concebida sem pecado, concebe o Filho pelo Espírito, tornando-se Sede da Sabedoria, nova Eva, Mãe da Igreja.

Pelo Espírito, manifesta o Filho aos pobres: pastores, magos, Simeão, Ana; torna-se Mulher na Cruz, orante em Pentecostes, inaugurando tempos finais.

Para catecúmenos adultos não iniciados, meditar João e Maria no rito de eleição quaresmal desperta desejo de pureza joanina e fiat mariano, preparando para o Batismo que unge com o Espírito na Vigília Pascal.

\section{Conclusão}

Professar o Espírito Santo revela o sopro vital da Trindade na fé e salvação. Adultos no RICA, abram-se ao Paráclito no catecumenato, para viver a unidade trinitária na iniciação